% sup1_03_kernel: Appendix on Ba3 Prop. 1, part 1, step S1: the kernel (1.16) of Ba3 as an external field (second alternative of (2.5)).

\subsubsection{S1: the kernel \cite[(1.16)]{Ba3} as an external field}\label{sup1:ssec:S1}

\cite[p.~414]{Ba3} expands $G_k(\Omega,\tilde A+\tilde B)$ by \cite[(3.44)--(3.45)]{Ba1} and treats the last term,
\begin{equation}\label{sup1:eq:116}
K_{n,n'}\defeq[G_k(\Omega,\tilde B)V_k(\tilde A,\tilde B)]^n\,G_k(\Omega,\tilde A+\tilde B)\,[V_k(\tilde A,\tilde B)G_k(\Omega,\tilde B)]^{n'},
\end{equation}
as an external field: ``for $n,n'$ sufficiently large, a kernel of the operator (1.16) is a sufficiently regular function of both variables \dots\ uniformly bounded by $O(1)(e(L^k\eps)^{1-\alpha})^{n+n'}$''. The consumed form is the second alternative of \cite[(2.5)]{Ba3}: $\norm{hK_{n,n'}h'}_{1,\alpha}\le O((e(L^k\eps)p(L^k\eps))^{n+n'})e^{-\delta_0\dist(\supp h,\supp h')}$ for the localisation functions $h,h'$ of the two end vertices. Here $h,h'$ are the smooth localisers of \cite[(1.33), p.~420]{Ba3} ($h=1$ on $\square(v)$, $h=0$ outside a neighbourhood, $|dh|,|ddh|\le c_h$), not the sharp cubes of the vertex localisations. At every use $m\defeq n+n'\ge\bar n+1\ge14$, since \cite[(3.45)]{Ba1} stops at $(n,n')=(\bar n+1,0)$ and $\bar n>12$ \cite[p.~421]{Ba3}; the proof needs only $m\ge6$ (or $m\ge4$ with $\alpha_0<\frac12$).

Write $G\defeq G_k(\Omega,\tilde B)$, $G'\defeq G_k(\Omega,\tilde A+\tilde B)$, $V\defeq G^{-1}-G'^{-1}=V_k(\tilde A,\tilde B)$ (the resolvent identity), so $K_{n,n'}=L_0VL_1V\cdots VL_m$ with $L_n=G'$ and $L_i=G$ for $i\ne n$.

\begin{lemma}[the vertex operator]\label{sup1:lem:S11}
Put $D\defeq D_{\tilde B}$, $\phi_\mu(x)\defeq\eta^{-1}(1-U(\tilde A_{x,\mu})^{-1})$, $\psi_\mu\defeq\phi_\mu^*\phi_\mu$, $\Pi\defeq P_k(\tilde A+\tilde B)-P_k(\tilde B)$, $\mathfrak D\defeq(1,D_1,D_2,D_3)$, $\supp\tilde A$ the set of endpoints of bonds with $\tilde A_b\ne0$, and $\mathcal U\defeq\{y:\dist_\infty(y,\supp\tilde A)\le1\}$.
\begin{enumerate}[label=\textup{(\roman*)},nosep,leftmargin=*]
\item $V=-[\sum_\mu(D_\mu^*\phi_\mu+\phi_\mu^*D_\mu+\psi_\mu)+a_k\Pi]=-\mathfrak D^*\mathfrak E\mathfrak D$, with $\mathfrak E_{00}=\sum_\mu\psi_\mu+a_k\Pi$, $\mathfrak E_{0\mu}=\phi_\mu^*$, $\mathfrak E_{\mu0}=\phi_\mu$, $0$ otherwise.
\item (Right-pushed form.) $V=-(R_0+\sum_\mu R_{1,\mu}D_\mu)=:-\mathfrak r\mathfrak D$, where $(R_{1,\mu}f)(x)\defeq\phi_\mu(x)^*f(x)-\phi_\mu(x-\eta e_\mu)U(\tilde B_{x-\eta e_\mu,\mu})^{-1}f(x-\eta e_\mu)$ and $R_0\defeq\sum_\mu(\psi_\mu-\mathsf m_\mu)+a_k\Pi$, $\mathsf m_\mu$ the multiplication by $(d^\eta_\mu\phi_\mu)(x-\eta e_\mu)$.
\item (Left-pushed form.) $V=V^*=-\mathfrak D^*\mathfrak r^*$.
\item Let $\sigma\le1$. Then $|\phi_\mu|,|d_\mu\phi_\mu|\le\sigma$, $|\psi_\mu|\le\sigma^2$, $|\Pi(y,y')|\le2d\sigma$ and $\sum_{y'}\eta^d|\Pi(y,y')|\le2d\sigma$; hence $\norm{\mathfrak r},\norm{\mathfrak E}\le c_V\sigma$ on $\ell^p(\eta^d)$, $p\in\{1,2,\infty\}$, $c_V\defeq4d+2da+1$. The kernels have sup-range $\le1$, and $V=\mathbf 1_{\mathcal U}V\mathbf 1_{\mathcal U}$, $\mathfrak r=\mathbf 1_{\mathcal U}\mathfrak r\mathbf 1_{\mathcal U}$, $\mathfrak E=\mathbf 1_{\mathcal U}\mathfrak E\mathbf 1_{\mathcal U}$.
\end{enumerate}
On $X=\Omega$ with Neumann conditions, (i)--(iv) hold if $\dist_\infty(\supp\tilde A,\Omega^c)\ge2$.
\end{lemma}

\begin{proof}
$U(\tilde A+\tilde B)=U(\tilde A)U(\tilde B)$, so $D_{\tilde A+\tilde B,\mu}f=U(\tilde A)(D_\mu f+\phi_\mu f)$, and since $U$ is pointwise orthogonal, $-\Delta_{\tilde A+\tilde B}=\sum_\mu(D_\mu+\phi_\mu)^*(D_\mu+\phi_\mu)$ on the same bond set; the masses cancel and the pins differ by $\Pi$: this is (i). (ii): $D_\mu^*(\phi_\mu f)(x)=-\phi_\mu(x-\eta e_\mu)U^{-1}_{x-\eta e_\mu,\mu}(D_\mu f)(x-\eta e_\mu)-(d_\mu\phi_\mu)(x-\eta e_\mu)f(x)$ (the computation of \cite[(2.32)]{Ba4}, redone in \cref{sup1:lem:S22}). (iii): both inverses are real symmetric. (iv): $|e^{q\theta}-e^{q\theta'}|\le|\theta-\theta'|$ gives $|\phi|\le e_k|\tilde A|$, $|d_\mu\phi_\mu|\le e_k|d\tilde A|$; for $y,y'$ in one block $B(z)$, $\Pi(y,y')=U(\tilde B\cdots)^T[U(\tilde A(\Gamma_{z,y}))^TU(\tilde A(\Gamma_{z,y'}))-1]U(\tilde B\cdots)$ with contours of length $\le d$, so $|\Pi|\le2de_k\norm{\tilde A}$; the Laplacian part lives on bonds with $\tilde A_b\ne0$ plus one $\eta$-shift, $\Pi$ on unit blocks whose contours meet $\supp\tilde A$, of sup-diameter $1-\eta$. On $\Omega$ every such bond and block lies in $\Omega$ once the distance is $\ge2$.
\end{proof}

\begin{lemma}[convolution and smearing; $d=3$, sup metric]\label{sup1:lem:S12}
\textup{(a)} For $\kappa_1,\kappa_2\in(0,3)$ and $0<\delta'<\delta\le1$: $\sum_w\eta^dK^\delta_{\kappa_1}(z-w)K^\delta_{\kappa_2}(w)\le c(\kappa_1,\kappa_2)(1+(\delta-\delta')^{-3})K^{\delta'}_{\kappa_1+\kappa_2}(z)$. \textup{(b)} $K^\delta_\kappa\le c(1+(\delta-\delta'')^{-1})K^{\delta''}_{\kappa'}$ for $\kappa\ge\kappa'$, $\delta''<\delta$. \textup{(c)} If $R$ is a multiplication, a one-bond shift, or a kernel supported in single unit blocks, with $|R(z,z')|\le r_R$ and column sums $\le r_R$, then $(|R|K^\delta_\kappa(\cdot-y))(z)\le c(\kappa)e^{2\delta}r_RK^\delta_\kappa(z-y)$.
\end{lemma}

\begin{proof}
(a) $e^{-\delta|z-w|-\delta|w|}\le e^{-\delta'|z|}e^{-(\delta-\delta')(|z-w|+|w|)}$. Put $r\defeq\max(|z|_\infty,\eta)$ and split into $|w|\le r/2$, $|z-w|\le r/2$ and the rest. In the two inner regions one factor is $\le cr^{\kappa_i-3}$ and the other is summed with $\sum_{|w|\le R}\eta^d\max(|w|,\eta)^{\kappa-3}\le c(\kappa)R^\kappa$ ($R\ge\eta$); the outer region is $\le c\int_{r/2}^\infty u^{\kappa_1+\kappa_2-4}e^{-(\delta-\delta')u}du$, which gives the stated classes (with the conventions at $\kappa_1+\kappa_2\ge3$). (b) For $|z|\le1$ the larger exponent gives the smaller kernel; for $|z|>1$ the power is absorbed by $e^{-(\delta-\delta'')|z|}$. (c) For $|z-y|_\infty\ge2$ every $z'$ with $R(z,z')\ne0$ has $\max(|z'-y|,\eta)\ge\frac12\max(|z-y|,\eta)$; for $|z-y|_\infty<2$ the block sum of $K_\kappa(\cdot-y)$ is $\le c(\kappa)$ and $K_\kappa(z-y)\ge ce^{-2\delta}$.
\end{proof}

\begin{lemma}[end factors]\label{sup1:lem:S13}
Let $L_0,L_1,L_2\in\{G,G'\}$, $a\in\{0,1\}$, $E^{(2)}_a\defeq D^aL_0VL_1$, $E^{(3)}_a\defeq D^aL_0VL_1VL_2$, every $V$ in right-pushed form; $\delta'\defeq\delta_K/2$, $0\le\beta\le\delta_K/4$. Assume (K) at every pair of points used (automatic on $T_\eta$).
\begin{enumerate}[label=\textup{(\roman*)},nosep,leftmargin=*]
\item For $y\in\mathcal U$: $\norm{E^{(2)}_a(x,y)}\le c_E\sigma K^{\delta'}_2(x-y)$ and $\norm{E^{(3)}_a(x,y)}\le c_E'\sigma^2K^{\delta'}_3(x-y)$.
\item For $a=1$ and $|x-\tilde x|\le1$ the $\tilde B$-transported difference in $x$ of $E^{(2)}_1$ is $\le c_E(\alpha)\sigma|x-\tilde x|^\alpha(K^{\delta'}_{2-\alpha}(x-y)+K^{\delta'}_{2-\alpha}(\tilde x-y))$; the same for $E^{(3)}_1$ with $\sigma^2$ and $K_{3-\alpha}$.
\item Put $J_\kappa\defeq(\sum_y\eta^dK^{\delta'}_\kappa(y)^2e^{2\beta|y|_\infty})^{1/2}$. Then $\norm{e^{\beta|\cdot-x|}\mathbf 1_{\mathcal U}E^{(2)}_a(x,\cdot)}_2\le c_EJ_2\sigma$; the transported difference of (ii) has weighted $\ell^2$ norm $\le2e^\beta c_E(\alpha)J_{2-\alpha}\sigma|x-\tilde x|^\alpha$; for $E^{(3)}$ the same with $J_3$, $J_{3-\alpha}$, $\sigma^2$. $J_\kappa<\infty$ uniformly in $\eta$ and in the torus side iff $\kappa>\frac32$: so $J_{2-\alpha}<\infty$ iff $\alpha<\frac12$, while $J_3$, $J_{3-\alpha}$ are finite for every $\alpha<1$.
\end{enumerate}
\end{lemma}

\begin{proof}
(i) $E^{(2)}_a=-(D^aL_0)R_0L_1-\sum_\mu(D^aL_0)R_{1,\mu}(D_\mu L_1)$: every line carries at most one derivative ($L_0$ the external one, $L_1$ the vertex one). For a line of type $G'$ write $D_{\tilde B,\mu}=U(\tilde A)^{-1}D_{\tilde A+\tilde B,\mu}-\phi_\mu$, so (K) for $G'$ gives $\norm{D_{\tilde B}G'(x,w)}\le c_K(K_1+\sigma K_2)\le cK^{\delta'}_1$ by \cref{sup1:lem:S12}(b). \Cref{sup1:lem:S12}(c) moves $R$ ($r_R\le c_V\sigma$) onto the kernels and (a), (b) give $K_{2-a}*K_1\to K_{3-a}\le cK_2$, $K_{2-a}*K_2\to K_2$; for $E^{(3)}_a=(D^aL_0)\mathfrak r(\mathfrak DL_1)\mathfrak r(\mathfrak DL_2)$ one gets $K_{2-a}*K_1*K_1\to K_3$. (ii) The difference acts on the first factor only. For $L_0=G$ use (K2), then $K_{1-\alpha}*K_1\to K_{2-\alpha}$ (resp.\ $K_{1-\alpha}*K_1*K_1\to K_{3-\alpha}$). For $L_0=G'$ write $U(\tilde B(\Gamma))=U((\tilde A+\tilde B)(\Gamma))U(\tilde A(\Gamma))^{-1}$: the quotient is at most the (K2) quotient of $D_{\tilde A+\tilde B}G'$, plus $3\sigma c_KK_1$ from $|U(\tilde A(\Gamma))^{-1}U(\tilde A_x)^{-1}-U(\tilde A_{\tilde x})^{-1}|\le e_k(\norm{\tilde A}|\Gamma|+\eta|\tilde A_x-\tilde A_{\tilde x}|)\le3\sigma|x-\tilde x|$, plus the $\phi$-terms, telescoped along $\Gamma$ as in \cref{sup1:lem:S15}, which give $c\sigma|x-\tilde x|(K_1(x-\cdot)+K_1(\tilde x-\cdot))$; since $|x-\tilde x|\le|x-\tilde x|^\alpha$ and $K_1\le cK_{1-\alpha}$, (K2) holds for $D_{\tilde B}G'$ with constant $c_K(\alpha)+10\sigma c_K$. (iii) Square and sum, with $|y-x|\le|y-\tilde x|+1$ for the $\tilde x$-term. Near $0$, $\sum_y\eta^d\max(|y|,\eta)^{2\kappa-6}$ is bounded uniformly in $\eta$ iff $2\kappa-6>-3$; the tail converges because $\delta'-\beta\ge\delta_K/4$.
\end{proof}

\begin{proposition}[the $\ell^2$ sandwich]\label{sup1:prop:S14}
Assume \eqref{sup1:eq:E}; then (H0), (K) for $G$ and $G'$, and $\sigma\le1$ hold. Let $0<\beta\le\beta_U\defeq\min\{\beta_*,\delta_K/4\}$, and either $m\ge4$ and $0<\alpha<\frac12$, or $m\ge6$ and $0<\alpha<1$. Put $\mathfrak N_{\alpha,\beta}(Y)\defeq\max_{a,b}\sup_{x,x'}e^{\beta|x-x'|_\infty}\max\{\norm{Y_{ab}(x,x')},aH_x,bH_{x'},abH_{xx'}\}$, $H$ the $\alpha$-H\"older quotients of $Y_{ab}$ with $\tilde B$-transport.
\textup{(a)} On $X=T_\eta$: $\mathfrak N_{\alpha,\beta}(K_{n,n'})\le C_U(\alpha)(C_V\sigma)^m$, with $C_V\defeq e^{2\beta}c_V\max\{c_W',1\}$, $c_W'\defeq(1+\sigma)^2c_W$.
\textup{(b)} On $X=\Omega$ with Neumann conditions and $\dist_\infty(\supp\tilde A,\Omega^c)\ge R_0+2$ the same bound holds with the supremum restricted to $x,x'$ such that $x$, $x'$, their bond partners, their H\"older partners and the points of the contours $\Gamma_{\tilde x,x}$, $\Gamma_{\tilde x',x'}$ have depth $\ge R_0$.
$C_U(\alpha)$ depends only on $d,a,M,L,\gamma_0,K_W,c_K,c_K(\alpha),\alpha,\beta$; the bound is uniform in $\eta$ (hence $k$, $\eps$), in the torus side, in $0\le m^2t^2\le c_m$ and in the backgrounds within \eqref{sup1:eq:E}.
\end{proposition}

\begin{proof}
$K_{n,n'}$ has $m$ vertices and $m+1$ lines, and each vertex is first order on each side (\cref{sup1:lem:S11}(i)), so $D^aKD^{b*}$ carries $m+2$ derivatives on $m+1$ lines and some line carries two; that line is put in the middle factor, which enters only through \cref{sup1:lem:W2}. No pointwise bound on $DGD^*$ is used.

\emph{Two-line ends ($m\ge4$).} $Y\defeq D^aKD^{b*}=E_L\mathcal NE_R$ with $E_L\defeq E^{(2)}_a=D^aL_0VL_1$, $\mathcal N\defeq VL_2V\cdots VL_{m-2}V$ (vertices $2,\dots,m-1$) and $E_R\defeq L_{m-1}VL_mD^{b*}=(D^bL_mVL_{m-1})^*$ (by $V=V^*$ and the symmetry of the $L_i$). Put the first vertex of $\mathcal N$ in right-pushed, the last in left-pushed and the others in symmetric form: $\mathcal N=(-1)^{m-2}\mathfrak r(\mathfrak DL_2\mathfrak D^*)\mathfrak E\cdots\mathfrak E(\mathfrak DL_{m-2}\mathfrak D^*)\mathfrak r^*$. By \cref{sup1:lem:S11}(iv) $\mathcal N=\mathbf 1_{\mathcal U}\mathcal N\mathbf 1_{\mathcal U}$, so $E_L(x,\cdot)$, $E_R(\cdot,x')$ are needed only on $\mathcal U$. For $L=G'$ each $\mathfrak D_{\tilde B}=U(\tilde A)^{-1}\mathfrak D_{\tilde A+\tilde B}-\Phi$ costs $(1+\sigma)$ per side; the kernels of $\mathfrak r,\mathfrak E$ have sup-range $\le1+\eta\le2$, so their conjugates by $e^{\pm\beta\rho}$ are bounded by $e^{2\beta}c_V\sigma$; \cref{sup1:lem:W2} bounds the blocks $\mathfrak DL\mathfrak D^*$. Hence $\norm{e^{\mp\beta\rho}\mathcal Ne^{\pm\beta\rho}}\le(e^{2\beta}c_V\sigma)^{m-2}c_W'^{m-3}$. Fix $x,x'$, $\rho\defeq|\cdot-x|_\infty$, and write $Y(x,x')=\ip{e^{\beta\rho}\mathbf 1_{\mathcal U}E_L(x,\cdot)^*}{(e^{-\beta\rho}\mathcal Ne^{\beta\rho})e^{-\beta\rho}\mathbf 1_{\mathcal U}E_R(\cdot,x')}$; Cauchy--Schwarz, $e^{-\beta|y-x|}\le e^{-\beta|x-x'|}e^{\beta|y-x'|}$ and \cref{sup1:lem:S13}(iii) (for $E_R$ by transposition) give $e^{\beta|x-x'|}\norm{Y(x,x')}\le(c_EJ_2\sigma)^2(e^{2\beta}c_V\sigma)^{m-2}c_W'^{m-3}$. $Y$ is linear in $E_L(x,\cdot)$ and in $E_R(\cdot,x')$ separately, so $H_x$ is the same pairing with $E_L(x,\cdot)$ replaced by $U(\tilde B(\Gamma_{\tilde x,x}))E_L(x,\cdot)-E_L(\tilde x,\cdot)$, similarly $H_{x'}$, $H_{xx'}$; \cref{sup1:lem:S13}(iii) bounds these with $2e^\beta c_E(\alpha)J_{2-\alpha}$, finite iff $\alpha<\frac12$.

\emph{Three-line ends ($m\ge6$, any $\alpha<1$).} Repeat with $E_L\defeq E^{(3)}_a$, $\mathcal N\defeq VL_3V\cdots VL_{m-3}V$ (vertices $3,\dots,m-2$, at least two since $m\ge6$) and $E_R\defeq L_{m-2}VL_{m-1}VL_mD^{b*}$; now $J_3$, $J_{3-\alpha}$ are finite. Since $c_V\ge1$ both cases give $\le C_U(\alpha)(C_V\sigma)^m$: (a).

(b) (H0) and \cref{sup1:lem:W2} hold on $\Omega$ with no depth condition, and \cref{sup1:lem:S11} holds since $\dist\ge2$. Every pointwise use of (K) above is at a point of $\mathcal U$ or a bond partner of one (depth $\ge R_0+2-1-\eta\ge R_0$), or at $x,x',\tilde x,\tilde x'$, their bond partners and the contour points of the telescoping in \cref{sup1:lem:S13}(ii), which are deep by the restriction. So (K) applies and the argument is verbatim.
\end{proof}

\begin{lemma}[localisation of two-variable bounds]\label{sup1:lem:S15}
Let $Y$ be a kernel, $\beta\ge0$, $M_Y\ge0$, and $\mathcal P,\mathcal P'$ sets of lattice points such that $e^{\beta|y-y'|_\infty}\max\{\norm{Y_{ab}(y,y')},aH_y,bH_{y'},abH_{yy'}\}\le M_Y$ for all $a,b$, all directions, all $y,\tilde y\in\mathcal P$, $y',\tilde y'\in\mathcal P'$ with partners $0<|y-\tilde y|,|y'-\tilde y'|\le1$. Let $0\le h,h'\le1$ with $|d_\mu h|,|d_\nu d_\mu h|\le c_h$ (and the same for $h'$), and assume that every point at sup-distance $\le1+\eta$ from $\supp h$ lies in $\mathcal P$, and likewise for $h'$, $\mathcal P'$. Then $\norm{hYh'}_{1,\alpha}\le e^{3\beta}(1+3\sqrt3c_h)^2M_Ye^{-\beta\dist_\infty(\supp h,\supp h')}$.
\end{lemma}

\begin{proof}
Leibniz: $D_\mu(hf)=F_1D_\mu f+F_2f$ with $F_1\defeq h(\cdot+\eta e_\mu)$, $F_2\defeq d_\mu h$; for the second variable, $(ZD_\mu^*)(x,x')=((D_\mu Z^\dagger)(x',x))^*$. Along a shortest monotone contour $\Gamma_{\tilde x,x}$, of $\ell^1$-length $\le\sqrt3|x-\tilde x|$, every $\varphi$ with $|d\varphi|\le c$ obeys $|\varphi(x)-\varphi(\tilde x)|\le\sqrt3c|x-\tilde x|$, and $U(\Gamma)W(x)-W(\tilde x)$ is a sum over the bonds of $\Gamma$ of transported $\eta D_\nu W$ terms, so $\norm{U(\Gamma)W(x)-W(\tilde x)}\le|x-\tilde x|_1\max_{z\in\Gamma,\nu}\norm{D_\nu W(z)}$. In one variable, $U(hY)_1(x)-(hY)_1(\tilde x)$ splits into $F_1(x)[UY_1(x)-Y_1(\tilde x)]$, $(F_1(x)-F_1(\tilde x))Y_1(\tilde x)$, $F_2(x)[UY_0(x)-Y_0(\tilde x)]$ (telescoped, $D_\nu Y_0$ being a component of $Y_1$) and $(F_2(x)-F_2(\tilde x))Y_0(\tilde x)$; with $|x-\tilde x|\le|x-\tilde x|^\alpha$ the quotient of $(hY)_1$ is $\le(1+3\sqrt3c_h)$ times the largest component or quotient of $Y$ at points of $\{x,\tilde x\}\cup\Gamma$, and a term vanishes unless $x$ or $\tilde x$ is within $\eta$ of $\supp h$, so every evaluated point is within $1+\eta$ of $\supp h$ and lies in $\mathcal P$. The expansions in $x$ and $x'$ commute (different variables, transports from opposite sides); the factors are $(1+c_h)^2$ per component, $(1+3\sqrt3c_h)(1+c_h)$ per one-variable quotient and $(1+3\sqrt3c_h)^2$ for the mixed one. Finally $|y-y'|_\infty\ge\dist_\infty(\supp h,\supp h')-2(1+\eta)$.
\end{proof}

\begin{corollary}[{\cite[(2.5)]{Ba3} for \cite[(1.16)]{Ba3} at $L=5$}]\label{sup1:cor:S16}
Assume \eqref{sup1:eq:E} and $\dist_\infty(\supp\tilde A,\Omega^c)\ge R_0+2$; let $m=n+n'\ge6$, $\alpha_0\in(0,1)$, $0<\beta\le\beta_U$, and let $h,h'$ be as in \cref{sup1:lem:S15} with supports at depth $\ge R_0+3$ in $\Omega$ (on $T_\eta$: no condition). Then
\[
\norm{hK_{n,n'}h'}_{1,\alpha_0}\le e^{3\beta}(1+6c_h)^2C_U(\alpha_0)\,(2c_1C_Ve_kp(t))^{n+n'}\,e^{-\beta\dist_\infty(\supp h,\supp h')}.
\]
\end{corollary}

\begin{proof}
Apply \cref{sup1:lem:S15} with $\mathcal P=\mathcal P'\defeq\{z:\rho_\Omega(z)\ge R_0+1+\eta\}$: for points of $\mathcal P$ the bond partners and the H\"older partners (and their bond partners, and the contour points between them) have depth $\ge R_0$, so \cref{sup1:prop:S14}(b) gives the hypothesis with $M_Y=C_U(C_V\sigma)^m$. A point within $1+\eta$ of a set of depth $\ge R_0+3$ has depth $\ge R_0+2-\eta\ge R_0+1+\eta$ since $\eta\le\frac15$. Finally $\sigma\le2c_1e_kp(t)$ and $3\sqrt3\le6$.
\end{proof}

\begin{remark}\label{sup1:rem:S1}
(i) This is the second alternative of \cite[(2.5)]{Ba3} with $\delta_0\defeq\beta\le\beta_U=\min\{\beta_*,\delta_K/4\}$; shrinking the rate is free. Since \cite[p.~426]{Ba3} extracts half of each line's rate for $\exp[-\frac12\delta_1\dist(\square(v),\square(v'))]$, the rate $\delta_0$ of Proposition~1 is $\min\{\delta_1/2,\beta\}>0$. (ii) The corollary holds for every $\alpha_0<1$, in particular $\alpha_0=\frac12$; for $m\in\{4,5\}$ it holds for $\alpha_0<\frac12$. It implies the \cite[p.~414]{Ba3} form read as $(e(L^k\eps))^{1-\alpha'}$ raised to $n+n'$: $(2c_1C_Ve_kp(t))^m\le(e_k^{1-\alpha'})^m$ as soon as $2c_1C_Vp(t)\le e_k^{-\alpha'}$, i.e.\ for $t\le t_0(\alpha')$, since $p$ is polylogarithmic and $e_k=et^{1/2}$. (The other parsing, $e\,(L^k\eps)^{1-\alpha}$, is not consumed: \cite[(1.33)]{Ba3} counts powers of $e(L^k\eps)$.) (iii) ``$n,n'$ sufficiently large'' is necessary: at $m=0$, with $g\defeq\eta^{-d}\mathbf 1_x$ and $\Lambda\defeq4d\eta^{-2}+m^2t^2+a\ge\norm H$, $(D_\mu G'D_\mu^*)(x,x)=\ip{D_\mu^*g}{G'D_\mu^*g}\ge\Lambda^{-1}\norm{D_\mu^*g}^2=2\eta^{-2-d}/\Lambda\sim\eta^{-3}/6$, not bounded uniformly in $\eta$. (iv) The depths used are $\supp\tilde A$ at $\ge R_0+2$ and vertex supports at $\ge R_0+3$, both in the sup metric and uniform in $k\ge1$.
\end{remark}
